\section{The reduced square root and a frozen matrix algebra}\label{sec:matrices}

The hierarchy of Proposition~\ref{prop:hierarchy} contains tensors of
arbitrarily high order whose derivatives have almost the same energy as the
tensors themselves.  We first transfer this near equality to the positive
square root of their one-slot reduction.  We then describe that square root
along the stationary diffusion by matrices frozen at the initial point.
The construction uses weighted traces throughout; in particular, no lower
bound on a nonzero eigenvalue will be required.

\subsection{Sharp energy for the reduced square root}

All limits in this section are along the counterexample sequence.  We make
the following diagonal choice explicit.  For each finite integer $K$, impose
the conclusions of Proposition~\ref{prop:hierarchy} through level $3K$, the
fixed-projection conclusions for rank coefficients at most $K$, and the
estimates of Lemma~\ref{lem:hermite-expansion} for every starting level
$0\le l<2K$ and every truncation order $0\le m\le K$, on time
intervals $0\le t\le K$.  Only levels through $3K$ enter these
expansions, including their remainders.  For each such finite list, the
hierarchy defects and the excesses over the normalized Hermite remainder
bounds can be made arbitrarily small.  Choose an integer $k\to\infty$
slowly enough that these defects and excesses tend to zero, uniformly over
the corresponding list.  Further auxiliary parameters below will be chosen
by the same finite-list procedure.  In particular, normalized endpoint
energies remain bounded when a band of $k$ levels is shifted by a fixed
number of levels.

Regard the direct sum
\[
 \Psi=\bigoplus_{k\le l<2k}\frac{T_l}{\sqrt{kI}}
\]
as a matrix by separating the first tensor index from the remaining
indices and the direct-sum label.  Define
\begin{equation}\label{eq:matrix-reduction}
 \rho=\Psi\Psi^*,\qquad R=\rho^{1/2}.
\end{equation}
High tensor degree provides many interchangeable slots, while averaging
adjacent degrees makes each fixed shift cost only endpoint terms of
order $k^{-1}$.  These two features will force sharp energy estimates
for the positive matrix field $R$.

Here $R$ is a positive matrix field, distinct from the scalar log-Sobolev
constant $R(\nu)$.  Write $\|\cdot\|_2$ for the field norm formed with
ordinary Hilbert--Schmidt coefficient norms.  The hierarchy gives
\begin{equation}\label{eq:matrix-normalization}
 \E\tr R^2=\|\Psi\|_2^2\longrightarrow1,
 \qquad \|D\Psi\|_2^2\longrightarrow1.
\end{equation}

\begin{lemma}[Centering and the square-root energy]\label{lem:matrix-centering}
The field $R$ belongs to the Dirichlet form domain and satisfies
\[
 \|\E R\|_{\HS}=o(1),\qquad
 \liminf\|DR\|_2^2\ge1,\qquad \|DR\|_2^2=O(1).
\]
\end{lemma}
\begin{proof}
For rectangular matrices $X,Y$, the Araki--Yamagami inequality
\cite[Theorem~1]{ArakiYamagami1981} gives the following square-root bound.
We give a finite-dimensional proof by self-adjoint dilation:
\begin{equation}\label{eq:matrix-root-lipschitz}
 \|\sqrt{XX^*}-\sqrt{YY^*}\|_{\HS}^2
 \le 2\|X-Y\|_{\HS}^2.
\end{equation}
Indeed the absolute value of the dilation of $X$ has diagonal blocks
$\sqrt{XX^*}$ and $\sqrt{X^*X}$.  If $A,B$ are self-adjoint and $f$ is
scalar Lipschitz, expansion in an eigenbasis of $A$ and an eigenbasis of $B$
shows that
$\|f(A)-f(B)\|_{\HS}\le\operatorname{Lip}(f)\|A-B\|_{\HS}$.
Apply this with $f(x)=|x|$.  Difference quotients, followed by form
approximation, prove that $R$ is in the form domain and that
$\|DR\|_2^2\le2\|D\Psi\|_2^2$.

The curvature conclusion of the hierarchy yields
\[
 \E\tr((1-M_s)R^2)=o(1).
\]
We also have
\begin{equation}\label{eq:matrix-independent-centering}
 \tr((\E M_s)(\E R^2))=o(1).
\end{equation}
To see the latter assertion in the strict branch, fix $\delta>0$ and let
$P_\delta$ be the spectral projection of $\E M_s$ onto $(\delta,\infty)$.
Lemma~\ref{lem:posterior-outliers} gives
$\rank P_\delta=O_\delta(I)$; the deterministic-projection conclusion of
Proposition~\ref{prop:hierarchy} therefore gives
$\tr(P_\delta\E R^2)=o(1)$.  Since $0\le M_s\le1$, the left side of
Equation~\eqref{eq:matrix-independent-centering} is at most
$o(1)+\delta\E\tr R^2$.  Send $\delta$ to zero after the sequence limit.
In the gap branch use directly $\|\E M_s\|_{\op}=o(1)$.

Let $R'$ be evaluated at an independent point.  Insert $1-M_s+M_s$
between $R$ and $R'$ in
$\|\E R\|_{\HS}^2=\E\tr(RR')$.
Hilbert--Schmidt Cauchy--Schwarz bounds the term containing $1-M_s$ by
\[
 \bigl(\E\tr((1-M_s)R^2)\bigr)^{1/2}
 \bigl(\E\tr R^2\bigr)^{1/2}=o(1).
\]
The other term is bounded by
\[
 \bigl(\E\tr(M_sR^2)\bigr)^{1/2}
 \bigl(\tr((\E M_s)(\E R^2))\bigr)^{1/2}=o(1).
\]
Poincar\'e for $\eta$, whose constant tends to one, now gives the lower
bound on $\|DR\|_2^2$ from Equation~\eqref{eq:matrix-normalization}.
\end{proof}

The next estimate is the reason for using tensors with many symmetric
slots.  It compares averaging an action over those slots with the
geometric, rather than arithmetic, weighting by the marginal matrix.

\begin{lemma}[Averaging over symmetric tensor slots]\label{lem:matrix-copied-slots}
Let $\psi$ be a finite direct sum of symmetric tensors, each with at least
$k\ge2$ slots.  Let $\rho$ be its one-slot reduction and $R=\sqrt\rho$.
For a self-adjoint matrix $H$, let $H_{\mathrm{av}}$ act on each component
by averaging the action of $H$ over its slots.  Then
\begin{equation}\label{eq:matrix-copied-slots}
 \|H_{\mathrm{av}}\psi\|^2
 \le\tr(RHRH)+\frac{2}{k}\tr(\rho H^2).
\end{equation}
\end{lemma}
\begin{proof}
For a component $\psi_d$ with $d$ slots, put
\[
 a_d=\langle\psi_d,H_{(1)}^2\psi_d\rangle,
 \qquad c_d=\langle\psi_d,H_{(1)}H_{(2)}\psi_d\rangle.
\]
Symmetry and Cauchy--Schwarz give $|c_d|\le a_d$, and direct expansion gives
\[
 \|H_{\mathrm{av}}\psi\|^2
 =\sum_d\left(\frac{a_d}{d}+\left(1-\frac1d\right)c_d\right)
 \le\sum_d c_d+\frac2k\sum_d a_d.
\]
It remains to prove $|\sum_d c_d|\le\tr(RHRH)$.  This correlation bound
can also be deduced from the Wigner--Yanase case of Lieb's trace-concavity
theorem \cite{Lieb1973}; we give a proof by weighted interpolation.

Define the linear map $C$ by
\[
 \tr(XC(Y))=\sum_d\langle\psi_d,X_{(1)}Y_{(2)}\psi_d\rangle.
\]
Flattening all remaining indices exhibits $C(Y)$ as a sum of terms
$A_\alpha Y^{\mathsf T}A_\alpha^*$.  Thus $C$ is positive and
$C(1)=C^*(1)=\rho$.  Restrict to the support of $\rho$, on which inverse
powers are defined, and set
\[
 U(Y)=\rho^{-1/2}C(Y)\rho^{-1/2}.
\]
The map $U$ is unital and contractive in operator norm: its preceding
representation is an isometric compression of copies of $Y^{\mathsf T}$.
It also satisfies $\tr(\rho U(Y))=\tr(\rho Y)$.
For $0\le\Re z\le1$ introduce
\[
 U_z(Y)=\rho^{z/2}
 U(\rho^{-z/2}Y\rho^{-z/2})\rho^{z/2}.
\]
On $\Re z=0$, multiplication on either side by imaginary powers of $\rho$
is isometric, so $U_z$ is operator-norm contractive.  At $z=1$,
\[
 U_1(Y)=C(\rho^{-1/2}Y\rho^{-1/2}),\qquad U_1^*(1)=1.
\]
Its adjoint is a unital positive map with the same compression and
transpose representation, hence is operator-norm contractive.  Duality
makes $U_1$ trace-norm contractive; imaginary powers give the same
conclusion on $\Re z=1$.

The interpolation used here is the finite-dimensional weighted Schatten
argument; compare \cite[Section~III, Theorem~4]{Beigi2013}.  We give the details for
our positive correlation map.
Given matrices of Hilbert--Schmidt norm one, continue the singular values
of the input by powers $2z$, and those of the matrix paired with its
output by powers $2(1-z)$, keeping the unitary factors fixed.  The bilinear
trace pairing is analytic in the strip, bounded by one on both boundary
lines by the preceding endpoint bounds and trace duality.  The three-lines
inequality bounds it by one at $z=1/2$.  Zero singular values follow by
approximation.  This proves $\|U_{1/2}(Y)\|_{\HS}\le\|Y\|_{\HS}$.
Take $Y=\rho^{1/4}H\rho^{1/4}$.  Then
\[
 \tr(YU_{1/2}(Y))=\tr(HC(H))=\sum_d c_d,
 \qquad \|Y\|_{\HS}^2=\tr(RHRH).
\]
Cauchy--Schwarz proves the claim.  Compression to the support leaves the
correlation unchanged: every tensor slot lies in that support.  Complexification, used
only for interpolation, preserves the conclusion for real tensors.
\end{proof}

\begin{lemma}[Saturation of the square-root energy]\label{lem:matrix-saturation}
There are measurable symmetric matrix fields $H_i$, $1\le i\le n$, solving
\begin{equation}\label{eq:matrix-sylvester}
 D_iR=\frac{H_iR+RH_i}{2},
\end{equation}
with zero kernel--kernel block, such that
\begin{equation}\label{eq:matrix-saturation}
 \|DR\|_2^2\longrightarrow1,\qquad
 \sum_i\E\|[H_i,R]\|_{\HS}^2=o(1),\qquad
 \sum_i\E\tr(R^2H_i^2)\longrightarrow1.
\end{equation}
Moreover,
\begin{equation}\label{eq:matrix-action-errors}
 \sum_i\left(\|D_iR-H_iR\|_2^2+
 \|D_i\Psi-(H_i)_{(1)}\Psi\|_2^2\right)=o(1).
\end{equation}
The first slot can be replaced by any existing slot in each component.
\end{lemma}
\begin{proof}
In an eigenbasis of $R$ with eigenvalues $r_a$, define
$(H_i)_{ab}=2(D_iR)_{ab}/(r_a+r_b)$ when $r_a+r_b>0$, and zero otherwise.
The kernel--kernel block of $D_iR$ vanishes almost everywhere.  One way to
check this for a form-domain positive matrix field is to use absolutely
continuous line representatives: at a differentiability point, the
quadratic form on every vector in the kernel has a local minimum zero,
and polarization gives the assertion.  The formula is measurable and gives
\[
 \sum_i\|H_iR\|_{\HS}^2\le4|DR|^2.
\]
Set
\[
 e=\|DR\|_2^2,\quad a=\sum_i\E\tr(R^2H_i^2),\quad
 c=\sum_i\E\|[H_i,R]\|_{\HS}^2,\quad d=\|D\Psi\|_2^2.
\]
The Sylvester equation implies the exact identities
\begin{equation}\label{eq:matrix-energy-identities}
 e=a-\frac c4,\qquad
 \sum_i\E\tr(RH_iRH_i)=e-\frac c4.
\end{equation}
Differentiating $R^2=\Psi\Psi^*$ in the weak form sense gives
\[
 e=\sum_i\E\langle D_i\Psi,(H_i)_{(1)}\Psi\rangle
   =\sum_i\E\langle D_i\Psi,(H_i)_{\mathrm{av}}\Psi\rangle.
\]
For example, the first equality follows by pairing $D_i(R^2)$ with $H_i/2$;
all products are integrable by the displayed bound.  The second uses the
symmetry of each $D_i\Psi$ on its existing slots.
Lemma~\ref{lem:matrix-copied-slots} and Cauchy--Schwarz yield
\[
 e^2\le d\left(e-\frac c4+\frac{2a}{k}\right).
\]
Since $a\le4e$, $d\to1$, and $\liminf e\ge1$, this first forces $e\to1$
and then $c\to0$; Equation~\eqref{eq:matrix-energy-identities} gives $a\to1$.
Finally,
\[
 \sum_i\|D_iR-H_iR\|_2^2=\frac c4,
 \qquad
 \sum_i\|D_i\Psi-(H_i)_{(1)}\Psi\|_2^2=d+a-2e.
\]
Both tend to zero.  Symmetry gives the assertion for the other slots.
\end{proof}

\subsection{Bounded columns and spectral blocks}

We next make the column $(H_i)_i$ bounded in operator norm.  This requires
a weighted tail estimate; an unweighted truncation would be inadequate
when $R$ has many small eigenvalues.  Throughout the modifications below,
$H_i$ denotes the current column; the exact Sylvester identity
\eqref{eq:matrix-sylvester} refers to the original column.  The approximate
derivative-action and balance estimates are preserved as stated in the
following lemmas.

\begin{lemma}[Balance and bounded columns]\label{lem:matrix-bounded-column}
The matrices in Lemma~\ref{lem:matrix-saturation} satisfy
\begin{equation}\label{eq:matrix-balance}
 \E\left\|\sum_iH_i\rho H_i-\rho\right\|_1=o(1),\qquad
 \E\left\|R\left(\sum_iH_i^2\right)R-\rho\right\|_1=o(1).
\end{equation}
They can be replaced by symmetric matrix fields, still denoted $H_i$, so
that Equation~\eqref{eq:matrix-action-errors}, the commutator conclusion of
Equation~\eqref{eq:matrix-saturation}, and
Equation~\eqref{eq:matrix-balance} hold, and
\begin{equation}\label{eq:matrix-column-bound}
 S:=\sum_iH_i^2\le4\Id,\qquad
 \E\|(S-1)R\|_{\HS}^2=o(1).
\end{equation}
The total change satisfies $\sum_i\E\|(H_i^{\mathrm{new}}-H_i^{\mathrm{old}})R\|_{\HS}^2=o(1)$.
\end{lemma}
\begin{proof}
We repeatedly use the dimension-free inequality
\begin{equation}\label{eq:matrix-reduction-perturbation}
 \|XX^*-YY^*\|_1
 \le(\|X\|_{\HS}+\|Y\|_{\HS})\|X-Y\|_{\HS}.
\end{equation}
By Equation~\eqref{eq:matrix-action-errors}, the first sum in
Equation~\eqref{eq:matrix-balance} is close in expected trace norm to the
one-slot reduction of $D\Psi$, with the derivative index included among
the column indices.  The derivative coupling identifies this with
\[
 \frac1{kI}\sum_{k\le l<2k}T_{l+1}T_{l+1}^*.
\]
The trace-norm error from the coupling is $o(1)$ by
Equation~\eqref{eq:matrix-reduction-perturbation} and
Equation~\eqref{eq:matrix-normalization}.  Shifting the band changes its
reduction only by its two endpoints, whose expected trace-norm cost is
\[
 \frac{\|T_{2k}\|_2^2+\|T_k\|_2^2}{kI}=O(k^{-1}).
\]
This proves the first balance equation.  Comparing the arrays of matrices
$(H_iR)_i$ and $(RH_i)_i$ proves the second by the same inequality and the
summed commutator estimate.

For the original, possibly unbounded column write
$\mathbf H v=(H_iv)_i$, $S=\mathbf H^*\mathbf H$, and $J=(1+S)^{-1}$.
Let $\mathbf K=\mathbf HR-R^{\oplus n}\mathbf H$, so that
$\E\|\mathbf K\|_{\HS}^2=o(1)$.  Algebra gives
\[
 [S,R]=\mathbf H^*\mathbf K-\mathbf K^*\mathbf H,
 \qquad [R,J]=J[S,R]J.
\]
The bounds
\[
 \|\mathbf HJ\mathbf H^*\|_{\op}\le1,\qquad
 \|\mathbf HJ\|_{\op}\le\tfrac12,\qquad \|J\|_{\op}\le1
\]
therefore imply $\|\mathbf H[R,J]\|_{\HS}\le2\|\mathbf K\|_{\HS}$.
For every fixed $f(x)=a+b/(1+x)$, it follows that
\[
 \E\|\mathbf H[f(S),R]\|_{\HS}^2=o(1).
\]
The squared norms of $\mathbf Hf(S)R$ and $\mathbf HRf(S)$ differ by $o(1)$:
the latter has bounded mean-square norm, and the preceding bound controls
their difference.  Testing the second balance equation against the bounded
matrix $f(S)^2$ consequently gives
\[
 \E\tr\bigl(R^2(S-1)f(S)^2\bigr)=o(1).
\]
Use $f=1,1+J,1-J$ and
$(S-1)^2/(1+S)=(S-1)-2(S-1)J$ to obtain
\begin{equation}\label{eq:matrix-column-tail}
 \E\tr\left(R^2\frac{(S-1)^2}{1+S}\right)=o(1).
\end{equation}
All these expressions are integrable: their absolute values are bounded
by a constant times $\tr(R^2(1+S))$.

Let $Q=\1_{(4,\infty)}(S)$ and $P=1-Q$.  On this spectral region,
$(S-1)^2/(1+S)$ dominates both a positive constant and a positive multiple
of $S$.  Thus
\[
 \E\tr(R^2Q)+\E\tr(R^2SQ)=o(1).
\]
The input truncation costs
$\sum_i\E\|H_iQR\|_{\HS}^2=\E\tr(R^2SQ)=o(1)$.
The output truncation costs
\[
 \sum_i\E\|QH_iR\|_{\HS}^2
 =\E\tr\left(Q\sum_iH_i\rho H_i\right)
 \le\E\tr(Q\rho)+\E\left\|\sum_iH_i\rho H_i-\rho\right\|_1=o(1).
\]
Consequently replacing $H_i$ by $PH_iP$ has total right-$R$ squared cost
$o(1)$.  The new column obeys
$\sum_i(PH_iP)^2\le PSP\le4P$.

Here is the perturbation accounting needed both now and below.  If
symmetric $\widetilde H_i$ have
\begin{equation}\label{eq:matrix-weighted-change}
 \Delta=\sum_i\E\|(\widetilde H_i-H_i)R\|_{\HS}^2=o(1),
\end{equation}
then symmetry gives the same bound with $R$ on the left.  Each derivative
error in Equation~\eqref{eq:matrix-action-errors} changes in norm by at most
$\sqrt\Delta$, and the summed commutator norm changes by at most
$2\sqrt\Delta$.  The two balance errors change in expected trace norm by
at most
\[
 \left(\Bigl[\sum_i\E\|H_iR\|_{\HS}^2\Bigr]^{1/2}
 +\Bigl[\sum_i\E\|\widetilde H_iR\|_{\HS}^2\Bigr]^{1/2}\right)\sqrt\Delta,
\]
by Equation~\eqref{eq:matrix-reduction-perturbation}, and this is $o(1)$.
The new summed energy is bounded by the old one plus the perturbation in
norm.  These statements apply to the compression just made.

For the bounded column, $[S,R]=\mathbf H^*\mathbf K-\mathbf K^*\mathbf H$
satisfies $\|[S,R]\|_2\le4\|\mathbf K\|_2=o(1)$.
Put $T=S-1$, so $\|T\|_{\op}\le3$.
The second balance equation says $\E\|RTR\|_1=o(1)$.  Finally,
\[
 \|TR\|_{\HS}^2
 =\tr(TRTR)+\tr(T[T,R]R).
\]
The expectation of the first term is at most $3\E\|RTR\|_1$ in absolute
value; the second is at most
$3\|[T,R]\|_2\|R\|_2=o(1)$.  This proves the last assertion.
\end{proof}

\begin{lemma}[Commutation and exchange of indices]\label{lem:matrix-index-exchange}
For the bounded columns just constructed,
\begin{equation}\label{eq:matrix-index-exchange}
 \sum_{i,j}\E\|[H_i,H_j]R\|_{\HS}^2=o(1),\qquad
 \E\sum_{i,a,b}|(H_iR)_{ab}-(H_aR)_{ib}|^2=o(1).
\end{equation}
These conclusions also hold after a symmetric perturbation satisfying
Equation~\eqref{eq:matrix-weighted-change}, provided the new column has a
uniform operator bound.
\end{lemma}
\begin{proof}
Equation~\eqref{eq:matrix-action-errors} implies that moving $H_i$ between
any two existing slots costs $o(1)$ in summed mean square.  Three slots
suffice to exchange two factors.  In the following chain, each $\simeq$
means that the difference has squared norm $o(1)$ after expectation and
summation over $i,j$:
\begin{align*}
 (H_i)_{(1)}(H_j)_{(1)}\Psi
 &\simeq (H_i)_{(1)}(H_j)_{(2)}\Psi
  = (H_j)_{(2)}(H_i)_{(1)}\Psi\\
 &\simeq (H_j)_{(2)}(H_i)_{(3)}\Psi
  = (H_i)_{(3)}(H_j)_{(2)}\Psi\\
 &\simeq (H_i)_{(3)}(H_j)_{(1)}\Psi
  = (H_j)_{(1)}(H_i)_{(3)}\Psi\\
 &\simeq (H_j)_{(1)}(H_i)_{(1)}\Psi.
\end{align*}
The equalities use commutation on different slots.  Each move costs at
most the one-factor error
multiplied by the column bound: for any tensor $v$,
\[
 \sum_i\|(H_i)_{(r)}v\|^2\le4\|v\|^2.
\]
The squared norm of the sum of the four errors is at most four times their
summed squared norms.  Thus the total commutator cost on $\Psi$ is $o(1)$.
Its cost on $\Psi$ equals its cost on $R$, since both have reduction $R^2$.

For the second conclusion, write $\Psi=RJ_0$, where $J_0J_0^*$ is the
support projection of $R$.  The matrix-valued array
$Z_{ia,b}=(H_iR)_{ab}-(H_aR)_{ib}$ has its column index in that support,
so $\|ZJ_0\|_{\HS}=\|Z\|_{\HS}$.  The entries of $ZJ_0$ express the
exchange of the derivative index and the first tensor index in
$(H_i)_{(1)}\Psi$.  Equation~\eqref{eq:matrix-action-errors} and the
approximation of $D\Psi$ by fully symmetric $T_{l+1}$ give the claim.
For perturbed matrices, the derivative errors and the slot-moving errors
have the bounds established after Equation~\eqref{eq:matrix-weighted-change};
the same four moves use the new uniform column bound.  The index-exchange
norm itself changes by at most $2\sqrt\Delta$.
\end{proof}

We have obtained a bounded column acting almost identically on the
symmetric tensor slots.  To use Gaussian matrix polynomials, we now put its
weighted estimates into a tracial form.  The following construction is
performed only at the initial point of the diffusion, so a varying
eigenbasis will never be differentiated.

\begin{lemma}[Spectral blocks and a bounded input map]\label{lem:matrix-blocks}
After enlarging the initial probability space by an independent uniform
variable, there are a positive block matrix
$R_*=\bigoplus_\beta\beta\Id_\beta$, with zero on $\ker R$, and symmetric
matrices $H_i$ supported on entries whose row, column, and outer index
belong to the same positive block, with the following properties.
Their total change satisfies Equation~\eqref{eq:matrix-weighted-change},
$R_*\le R\le(1+\epsilon)R_*$ on the positive support for
$\epsilon\to0$, and the column bound $S\le4\Id$ holds.  With the positive
weighted trace
\begin{equation}\label{eq:matrix-weighted-trace}
 \tau(U)=\E\sum_\beta\beta^2\tr_\beta U_\beta,
\end{equation}
where the expectation also includes any fresh noises present, we have
$\tau(1)\to1$ and
\begin{equation}\label{eq:matrix-tracial-defects}
 \tau((S-1)^2)=o(1),\quad
 \sum_{i,j}\tau(|[H_i,H_j]|^2)=o(1),\quad
 \E\sum_\beta\beta^2\sum_{i,a,b\in\beta}
 |(H_i)_{ab}-(H_a)_{ib}|^2=o(1).
\end{equation}
For a vector $v$ in a positive block define $H[v]=\sum_{i\in\beta}v_iH_i$.
We can further impose, pointwise,
\begin{equation}\label{eq:flattening-bound}
 \|H[v]\|_{\HS}\le2|v|.
\end{equation}
The bound uses ordinary, unnormalized Hilbert--Schmidt norm.
All the derivative-action and balance conclusions above continue to hold
for these matrices in the initial coordinates.
\end{lemma}
\begin{proof}
Choose an orthogonal eigenframe $O$ for $R$ at the initial point.  Transform
all three indices: in that frame the matrices are
\[
 \widehat H_i=\sum_jO_{ji}O^*H_jO.
\]
This preserves the column bound, all the summed norms, and the
index-exchange estimate, since these are Euclidean tensor contractions.
Finite-dimensional spectral calculus and an ordered orthonormal basis
completion give measurable choices, including at repeated eigenvalues.

Let $h=\log(1+\epsilon)$ and take an independent $\theta$ uniform on
$[0,h)$.  Group positive eigenvalues in
$[e^{\theta+mh},e^{\theta+(m+1)h})$, $m\in\mathbb Z$, and round those in
each nonempty block down to $\beta=e^{\theta+mh}$.  The zero eigenspace is
kept separate.  Write
\[
 \kappa=\sum_i\E\|[H_i,R]\|_{\HS}^2=o(1),\qquad
 a=\sum_i\E\|H_iR\|_{\HS}^2=O(1).
\]
For $0<\delta<h$, a pair of positive eigenvalues with
$|\log(r_a/r_b)|\ge\delta$ satisfies
$r_b^2\le C\delta^{-2}(r_a-r_b)^2$, when $\delta\le1$.
If their logarithmic distance is smaller than $\delta$, the probability
that a block boundary separates them is at most $\delta/h$.  Consequently
the right-$R$ squared cost of deleting entries with row and column in
different blocks is bounded by
\begin{equation}\label{eq:matrix-block-error}
 C\delta^{-2}\kappa+\frac{\delta}{h}a+\kappa.
\end{equation}
The final term covers a zero row and positive column directly by the
commutator.  A zero column has no right-$R$ cost.  For example, choose
$\delta\to0$ slowly enough that $\kappa/\delta^2\to0$, and then
$h\to0$ slowly enough that $\delta/h\to0$.

Deleting entries whose outer index belongs to another block has vanishing
cost as well.  On the set where $i$ and $b$ have different block labels,
\[
 |(H_i)_{ab}r_b|^2
 \le2|(H_i)_{ab}r_b-(H_a)_{ib}r_b|^2
   +2|(H_a)_{ib}r_b|^2.
\]
Sum over the indices.  The first term is the second defect of
Equation~\eqref{eq:matrix-index-exchange}; the second is a row--column block
exclusion of the kind already bounded in
Equation~\eqref{eq:matrix-block-error}.  This includes outer indices in the
zero eigenspace.  Within each block the resulting matrices have the form
$P_\beta\widehat H_iP_\beta$ for $i\in\beta$, and vanish otherwise.  Thus
\[
 \sum_{i\in\beta}(P_\beta\widehat H_iP_\beta)^2
 \le P_\beta\left(\sum_i\widehat H_i^2\right)P_\beta\le4P_\beta.
\]
They are symmetric and their total change has the cost in
Equation~\eqref{eq:matrix-weighted-change}.

For every block matrix $X$ one has the precise comparison
\begin{equation}\label{eq:matrix-weight-comparison}
 \tau(X^*X)\le\E\|XR\|_{\HS}^2
 \le(1+\epsilon)^2\tau(X^*X).
\end{equation}
In particular $\tau(1)\to1$.
The perturbation bounds following Equation~\eqref{eq:matrix-weighted-change}
prove the derivative, commutator, and balance conclusions for these
matrices.  The bounded-column argument at the end of
Lemma~\ref{lem:matrix-bounded-column} again gives
$\E\|(S-1)R\|_{\HS}^2=o(1)$, so
Equation~\eqref{eq:matrix-weight-comparison} proves the first tracial defect.
Lemma~\ref{lem:matrix-index-exchange} gives the second.  For the third,
all three indices now lie in one block, and the index-exchange defect
is exactly weighted by $r_b^2$; replacing this by $\beta^2$ only decreases
it.  Matrix symmetry and exchange of $i,a$ together give asymptotic
symmetry under every permutation of the three indices.

It remains to bound the map from the outer index to matrix coefficients.
On each block let
\[
 Fv=\sum_i v_iH_i,\qquad (Gv)_{ab}=(H_bv)_a.
\]
The map $G$ has operator norm at most two because
$\|Gv\|_{\HS}^2=\langle v,Sv\rangle\le4|v|^2$.
Total symmetry just established says
\[
 \E\sum_\beta\beta^2\|F-G\|_{\HS(\R^{\dim\beta},\HS)}^2=o(1).
\]
Clip the singular values of $F$ to two, obtaining $F'=FB$ for a real
positive contraction $B$ on the input space.  Metric projection onto the
operator-norm ball in Hilbert--Schmidt norm gives
$\|F-F'\|_{\HS}\le\|F-G\|_{\HS}$; this also follows immediately by
singular-value decomposition and the variational characterization of
singular values.  Thus the block-weighted squared cost is $o(1)$.
Writing $H'_i=\sum_j B_{ji}H_j$ shows that every $H'_i$ remains symmetric.
For each vector $v$,
\[
 \sum_i\|H'_iv\|^2\le\sum_j\|H_jv\|^2,
\]
so the column bound is preserved, as is the block support.  The weighted
cost controls the right-$R$ cost by
Equation~\eqref{eq:matrix-weight-comparison}.  Therefore the derivative and
balance errors have exactly the bounds stated after
Equation~\eqref{eq:matrix-weighted-change}; index exchange changes by at
most twice the perturbation norm, and the four-move argument controls
the commutators.  Repeating the bounded-column last step and then
Equation~\eqref{eq:matrix-weight-comparison} proves all three tracial defects
for $H'_i$.  The defining singular-value clipping proves
Equation~\eqref{eq:flattening-bound}.
\end{proof}

\subsection{Gaussian moments and the frozen diffusion law}

The matrices within a block need not commute exactly.  The next lemma
shows that their Gaussian linear combinations nevertheless have the
scalar Gaussian joint moments in the weighted limit.  We include the
moment bounds because they also justify the unbounded exponential needed
for the diffusion law.

\begin{lemma}[Weighted Gaussian moments]\label{lem:matrix-gaussian-moments}
Let the block matrices and $\tau$ be those of
Lemma~\ref{lem:matrix-blocks}.  For independent standard Gaussian vectors
$N_1,\ldots,N_q$, independent of the initial data, put $D_r=H[N_r]$.
Then, for every fixed word $r_1,\ldots,r_m$,
\begin{equation}\label{eq:matrix-gaussian-joint-law}
 \tau(D_{r_1}\cdots D_{r_m})
 \longrightarrow \E(Z_{r_1}\cdots Z_{r_m}),
\end{equation}
where $Z_1,\ldots,Z_q$ are independent scalar standard normals.
Products of any fixed length can be reordered with error tending to zero
in $L^2(\tau)$.  For every integer $p\ge1$ and real $u\ge0$,
\begin{equation}\label{eq:matrix-gaussian-moment-bound}
 \E_{N_r}\tr_\beta D_r^{2p}
 \le4^p(2p-1)!!\tr_\beta1,
 \qquad
 \tau(e^{u|D_r|})\le2e^{2u^2}\tau(1).
\end{equation}
If $\mathfrak h_j$ is the tensor Hermite polynomial of
Lemma~\ref{lem:hermite-expansion} and $\mathfrak h_j^{(1)}$ is its
one-dimensional version, then for each fixed $j$,
\begin{equation}\label{eq:matrix-wick-reduction}
 \tau\left(\left|
 \sum_{i_1,\ldots,i_j}H_{i_j}\cdots H_{i_1}
      \mathfrak h_j(N_1)_{i_1\ldots i_j}
 -\mathfrak h_j^{(1)}(D_1)\right|^2\right)=o(1).
\end{equation}
\end{lemma}
\begin{proof}
We first prove a conditional moment bound on one fixed block, allowing a
general column bound $S\le C\Id$.  We use the Gaussian trace-moment
integration-by-parts argument of \cite[Lemma~6.1 and Section~7]{Tropp2018};
the weighted near-commutation limit will be proved afterward.
Write $D=\sum_iN_iH_i$.
Gaussian integration by parts gives
\[
 \E_N\tr D^{2p}
 =\sum_{j=0}^{2p-2}\E_N\sum_i
   \tr(H_iD^jH_iD^{2p-2-j}).
\]
In an eigenbasis of $D$ with eigenvalues $d_a$, the absolute value of a
summand is bounded by
\[
 \sum_{i,a,b}|(H_i)_{ab}|^2|d_b|^j|d_a|^{2p-2-j}.
\]
For $p>1$, the scalar weighted arithmetic--geometric mean inequality
bounds the product of powers by
\[
 \frac{j}{2p-2}|d_b|^{2p-2}
 +\left(1-\frac{j}{2p-2}\right)|d_a|^{2p-2}.
\]
Since $\sum_{i,b}|(H_i)_{ab}|^2=S_{aa}\le C$ and similarly with $a,b$
exchanged, each summand is at most $C\tr|D|^{2p-2}$.
For $p=1$ the assertion is $\tr S\le C\tr1$.
Induction proves
$\E_N\tr D^{2p}\le C^p(2p-1)!!\tr1$.
After multiplying by the block weights, all fixed moments are uniformly
bounded.  Also $e^{u|x|}\le2\cosh(ux)$, so the even-moment series gives
\[
 \tau(e^{u|D|})
 \le2\sum_{p\ge0}\frac{u^{2p}C^p(2p-1)!!}{(2p)!}\tau(1)
 =2e^{Cu^2/2}\tau(1).
\]
This proves Equation~\eqref{eq:matrix-gaussian-moment-bound} with $C=4$.

For clarity, the norm $\|X\|_{p,\tau}=\tau(|X|^p)^{1/p}$ is formed on
the direct sum of the blocks and then averaged over the classical random
variables.  Schatten H\"older followed by scalar H\"older gives
\[
 \|X_1\cdots X_m\|_{p,\tau}
 \le\prod_{j=1}^m\|X_j\|_{p_j,\tau},
 \qquad \sum_j\frac1{p_j}=\frac1p.
\]
These bounds use the actual weights $\beta^2$, with no factor depending
on a block dimension.  The finite-matrix Schatten inequality can itself
be obtained by the same singular-value three-lines argument as in
Lemma~\ref{lem:matrix-copied-slots}, interpolating the endpoint bounds
with one trace norm and the other norms operator norms, and then applying
duality.

For independent $N_r,N_s$,
\[
 \|[D_r,D_s]\|_{2,\tau}^2
 =\sum_{i,j}\tau(|[H_i,H_j]|^2)=o(1).
\]
All its higher fixed moments are bounded by H\"older and the moment bound
just proved.  Interpolation between $L^2(\tau)$ and an arbitrarily high
fixed $L^p(\tau)$ therefore makes this commutator tend to zero in every
fixed finite power norm.  Applying H\"older to the factors on either side
of a commutator shows that an adjacent transposition in a word of fixed
length costs $o(1)$ in $L^2(\tau)$.  A finite sequence of transpositions
proves the reordering assertion.

Now expand a mixed moment by Gaussian pairings, conditional on the
matrices.  Pairings must join indices of the same Gaussian vector.  For
each pair introduce its own independent standard Gaussian vector, shared
only between the two factors in that pair.  This represents the
contracted matrix word exactly as a conditional expectation.  Reorder the
word until each pair is adjacent; the error is $o(1)$ even in
$L^2(\tau)$.  Averaging the independent vectors then gives $S^{m/2}$.
For each fixed integer $b$, boundedness of $S$ gives
\[
 \|S^b-1\|_{2,\tau}
 \le\left(\sum_{a=0}^{b-1}4^a\right)\|S-1\|_{2,\tau}=o(1).
\]
Each pairing consequently contributes $\tau(1)+o(1)=1+o(1)$.
Odd moments vanish.  The number of admissible pairings is precisely the
scalar independent-normal moment, proving
Equation~\eqref{eq:matrix-gaussian-joint-law}.

Finally expand the tensor Hermite polynomial using its defining
exponential
$\exp(\langle v,N\rangle-|v|^2/2)$.
A term consists of free $N$ indices and a number of signed Kronecker-delta
pairs.  Represent each delta pair by its own independent Gaussian vector,
as above, and move its two matrix factors together.  Conditional averaging
is contractive in $L^2(\tau)$, so the reordering error remains $o(1)$.
The resulting factors $S$ can be replaced by one, even between fixed
powers of $D_1$, by H\"older and interpolation applied to the bounded
matrix $S-1$.  The resulting polynomial is
\[
 \sum_{b=0}^{\lfloor j/2\rfloor}
 \frac{(-1)^b j!}{2^bb!(j-2b)!}D_1^{j-2b}
 =\mathfrak h_j^{(1)}(D_1).
\]
There are only finitely many terms for fixed $j$, proving
Equation~\eqref{eq:matrix-wick-reduction}.
\end{proof}

\begin{proposition}[Frozen diffusion law]\label{prop:frozen-law}
Let $Y_u$ be the stationary diffusion with invariant law $\eta$ and driving
Brownian motion $B_u$, and fix $t>0$.  Make the matrix construction of
Lemma~\ref{lem:matrix-blocks} at $Y_0$, using an independent random block
offset.  In its orthonormal frame put $N=B_t/\sqrt t$; conditionally on the
initial data and block offset, $N$ is a standard Gaussian vector.  Define,
in the initial spatial coordinates,
\begin{equation}\label{eq:matrix-frozen-covariance}
 \sigma=\bigoplus_\beta\beta^2
       \exp\bigl(\sqrt{8t}\,H[N]-4t\Id_\beta\bigr),
\end{equation}
and set $\sigma=0$ on the omitted kernel.  Then
\begin{equation}\label{eq:matrix-frozen-law}
 \E\|R(Y_t)-\sigma^{1/2}\|_{\HS}^2\longrightarrow0.
\end{equation}
The matrices defining $\sigma$ obey the bounded-column and input-map
bounds \eqref{eq:matrix-column-bound} and \eqref{eq:flattening-bound},
the weighted defects \eqref{eq:matrix-tracial-defects}, and the mixed
moment law of Lemma~\ref{lem:matrix-gaussian-moments}.
\end{proposition}
\begin{proof}
The proof has three steps: replace the added tensor indices in the
hierarchy expansion by products of the initial matrices, identify their
Wick contractions, and sum the resulting scalar Hermite series.
All limits involving the degree of that series are taken after the
counterexample-sequence limit, with $t$ fixed.

The derivative-action estimate and the hierarchy coupling give
\begin{equation}\label{eq:matrix-one-step}
 \frac1{kI}\sum_{k\le l<2k}\sum_i
 \|T_{l+1,i\,\cdot}-(H_i)_{(1)}T_l\|_2^2=o(1).
\end{equation}
It also holds when the band is shifted by any fixed integer $a\ge0$.
Indeed the same matrices have been constructed using levels $k$ through
$2k-1$.  Removing or adding at most $a$ levels in the error sum costs
$O_C(a/k)$, since the unnormalized squared norm of each one-step
error is at most
\[
 2\|T_{l+1}\|_2^2+2C\|T_l\|_2^2=O_C(I),
\]
by the bounded column and the diagonal endpoint control.  This also
explains why the matrices may already include the block and input
clipping modifications: their cost on $\Psi$ is the right-$R$ cost in
Equation~\eqref{eq:matrix-weighted-change}.

Iterating Equation~\eqref{eq:matrix-one-step}, for each fixed $j$ we may
replace the extra-index tensor $T_{l+j}$ by
\[
 (i_1,\ldots,i_j)\longmapsto
 (H_{i_j})_{(1)}\cdots(H_{i_1})_{(1)}T_l
\]
in the normalized band mean-square norm, with error $o(1)$.
More explicitly, if $e_a$ is the root mean-square error in the one-step
estimate on the band shifted by $a$, the root error after $j$ steps is at
most
\[
 \sum_{a=0}^{j-1}C^{a/2}e_{j-1-a}.
\]
Here multiplication by a column of $H_i$ costs at most $\sqrt C$ in the
summed coefficient norm.  Each $e_a=o(1)$ at fixed $a$, so the displayed
bound vanishes.  Contraction with $\mathfrak h_j(N)$ multiplies squared
coefficient norm by at most $j!$: it first symmetrizes the contracted
indices, an orthogonal projection, and Gaussian Hermite orthogonality
then gives the factor $j!$.

Apply Lemma~\ref{lem:hermite-expansion} to the band of levels in $\Psi$.
For each fixed $m$, the preceding replacements and
Equation~\eqref{eq:matrix-wick-reduction} give
\[
 \Psi(Y_t)=P_{m,t}(D)\Psi(Y_0)+\mathcal E_{m,t},\qquad D=H[N],
\]
where
\[
 P_{m,t}(x)=e^{-t}\sum_{j=0}^m
          \frac{(2t)^{j/2}}{j!}\mathfrak h_j^{(1)}(x)
\]
and
\begin{equation}\label{eq:matrix-hermite-error}
 \limsup\|\mathcal E_{m,t}\|_2
 \le \frac{(2t)^{(m+1)/2}}{\sqrt{(m+1)!}}.
\end{equation}
To justify the conversion of a tracial polynomial error to the norm of
its action on $\Psi(Y_0)$, use
$\|X\Psi(Y_0)\|_{\HS}=\|XR\|_{\HS}$ and
Equation~\eqref{eq:matrix-weight-comparison}.  The band version of the
remainder bound follows from the uniform finite-list diagonal fixed at
the start of the section.  Its right side tends to zero as $m\to\infty$.

The spectral measures of $D$ averaged with $\tau$ have total mass tending
to one, converge in moments to the standard Gaussian by
Equation~\eqref{eq:matrix-gaussian-joint-law}, and have uniformly bounded
exponential moments at every fixed linear parameter by
Equation~\eqref{eq:matrix-gaussian-moment-bound}.  These facts imply weak
convergence to the standard Gaussian and convergence of the integrals of
every continuous function bounded by a polynomial times a fixed linear
exponential.  Here is a direct justification: the second moment gives
tightness; the higher moment bounds give uniform integrability of each
polynomial, so every subsequential limit has the Gaussian moments.  Its
exponential moments determine its characteristic function by its power
series, hence determine the law.  Using a larger exponential parameter
makes the tails of any of the specified test functions uniformly small.

The scalar Hermite generating identity gives, with $Z$ standard normal,
\[
 e^{-t}\sum_{j\ge0}\frac{(2t)^{j/2}}{j!}\mathfrak h_j^{(1)}(Z)
 =\exp(\sqrt{2t}\,Z-2t)
\]
in squared mean.  In fact its squared tail norm is
$e^{-2t}\sum_{j>m}(2t)^j/j!$.  Define the matrix multiplier
\[
 M=\exp(\sqrt{2t}\,D-2t\Id).
\]
The just-proved convergence of spectral integrals, applied with each
fixed $m$, yields
\begin{equation}\label{eq:matrix-exponential-limit}
 \lim\tau(|M-P_{m,t}(D)|^2)
 =e^{-2t}\sum_{j>m}\frac{(2t)^j}{j!},\qquad
 \tau(M^2)\longrightarrow1.
\end{equation}
Together with Equations~\eqref{eq:matrix-hermite-error} and
\eqref{eq:matrix-weight-comparison}, first sending the sequence index and
then $m$ to infinity, this proves
\[
 \E\|\Psi(Y_t)-M\Psi(Y_0)\|_{\HS}^2=o(1).
\]
There is no operator-norm bound on $M$ in this argument; only its weighted
trace moments are used.

Write $\Psi(Y_0)=RJ_0$, with $J_0J_0^*$ equal to the support projection.
On each block $(R-R_*)^2\le\epsilon^2R_*^2$, and $R_*$ is scalar there.
By cyclicity and positivity of trace,
\[
 \E\|M(R-R_*)J_0\|_{\HS}^2
 \le\epsilon^2\tau(M^2)=o(1).
\]
Thus $\Psi(Y_t)$ is close in squared mean to $MR_*J_0$.
Its reduction is exactly
\[
 (MR_*J_0)(MR_*J_0)^*
 =MR_*^2M=\bigoplus_\beta\beta^2
       \exp(\sqrt{8t}\,H[N]-4t\Id_\beta)=\sigma.
\]
Finally Equation~\eqref{eq:matrix-root-lipschitz} transfers the approximation
of these rectangular matrices to the approximation of their positive
reduced square roots, proving Equation~\eqref{eq:matrix-frozen-law}.
\end{proof}

The resulting description retains the weights $\beta^2$ and the ordinary
trace on every block.  It supplies both a bounded matrix column and the
input bound of Equation~\eqref{eq:flattening-bound}, while allowing any
number of positive blocks and eigenvalues arbitrarily close to zero.
These are the properties needed for the covariance observations in
Section~\ref{sec:observations}.
